\documentclass{article}
\usepackage[letterpaper,margin=0.45in]{geometry}
\usepackage[T1]{fontenc}
\usepackage{newtxtext,newtxmath}
\usepackage{microtype}
\usepackage{array,tabularx,booktabs}
\usepackage{enumitem}
\pagestyle{empty}
\setlength{\parindent}{0pt}
\setlength{\parskip}{1.5pt}
\setlength{\tabcolsep}{2.8pt}
\renewcommand{\arraystretch}{1.04}
\newcolumntype{Y}{>{\raggedright\arraybackslash}X}
\begin{document}
\fontsize{8.6}{9.35}\selectfont

\begin{center}
{\bf Author Response: CHORD Hallucination Mitigation}
\end{center}

We thank the reviewers for identifying that aggregate benchmark scores alone do not isolate mechanism, detector attribution, and deployment cost. We therefore add targeted diagnostics and narrow the claim: CHORD is a detector-assisted, base-MLLM-training-free admission-time verifier for object-grounded hallucination mitigation; decoder-to-vision attention is used as an operational scoring signal, not as a causal explanation.

\vspace{1pt}
\begin{tabularx}{\linewidth}{@{}p{0.165\linewidth}Y Y Y@{}}
\toprule
Concern & Reviewer need & Response / evidence to add & Boundary in final paper \\
\midrule
Future mechanism & M8du asks whether Future actually changes admitted tokens and whether those changes help. & We add a paired Full-vs-P+C admission diagnostic: rollout-induced token changes, corrected unsupported mentions, and harmful changes. Exact flip/correctness counts are included only for completed paired logs; otherwise the rebuttal states the diagnostic protocol without expected numbers. & Future is a quality-oriented verifier, not a claim that every token needs rollout. If paired logs are unavailable, we do not report expected percentages, counts, confidence intervals, or p-values. \\
\midrule
Detector attribution & KrEs/M8du ask whether Grounding DINO alone explains the gains. & We isolate detector contribution with fixed same-anchor, uniform/random-anchor, no-Current, P+C, and Full controls. If control logs are complete, the compact attribution row is reported; otherwise we state the control design and limitation. & CHORD is detector-assisted. We do not claim detector independence or second-proposer robustness unless those runs are measured. \\
\midrule
Cost/defaults and baselines & jjVG/KrEs/ve3y ask about overhead, recent baselines, and why $k=5,m=3$. & We separate proposal cost from decode ITL. P+C is the practical/default regime; Full is quality-oriented/offline. $k=5,m=3$ is the quality setting, while P+C or $k=5,m=2$ is practical under latency constraints. ONLY, VHD/VHR, and HALC use matched backbone, split, prompt family, parser, seeds, and validation budget. & We report numeric total cost, VRAM, batch behavior, and recent-baseline wins only when measured with provenance; otherwise they remain camera-ready additions, not rebuttal evidence. \\
\bottomrule
\end{tabularx}

\vspace{2pt}
We will redraw Fig.~2 as sequential Past/Current/Future stages. The camera-ready version will remove detector-independent, broad relation/composition, or causal-attention claims unless supported by measured evidence. All exact numbers in the official page must be measured or removed; expected target values from the internal master are not used as factual rebuttal evidence.

\end{document}
